\section{Method}
\label{sec:proposed_solution}

This section presents the role-shared semantic coordination framework shown in
Figure~\ref{fig:method_overview}. Its three roles respectively evaluate scene
feasibility, decompose and assign task dependencies, and verify outcomes while
identifying event-affected subtasks. Section~\ref{sec:role_shared_coordination}
defines the shared semantic state and these three responsibilities,
Section~\ref{sec:semantic_physical_coordination} describes event-scoped
updates, Section~\ref{sec:dispatch_execute} constructs and verifies executable
commands, and Section~\ref{sec:closed_loop_coordination} gives the online
algorithm and complexity discussion.

\subsection{Construct the role-shared semantic state}
\label{sec:role_shared_coordination}

The shared semantic state is the authoritative task-level state maintained by
the three semantic roles during the current coordination cycle:
\begin{equation}
 S_\nu^{\texttt{sem}}\triangleq
 \bigl(\gamma,G,\mathcal{F},z\bigr),
 \label{eq:semantic_state}
\end{equation}
where $\nu$ is the semantic version, $\gamma$ is the current task instruction,
$G$ is the assignment-aware primitive graph,
$\mathcal{F}$ contains task-relevant scene and feasibility facts, and $z$
records primitive-subtask execution status. This state excludes
manipulator-local control variables and stores only the semantic information
used for task-level coordination.

The affordance role maintains $\mathcal{F}$ through
\begin{equation*}
 \rho_A:o(t)\mapsto\mathcal{F}.
\end{equation*}
It extracts task-relevant scene and
feasibility facts from the observation stream, including object grounding,
relations, affordances, and feasible primitive interactions. During
initialization, observed objects are grounded to semantic targets in
$\gamma$. After each primitive execution, the latest observation refreshes
$\mathcal{F}$ before planning or verification consumes the updated physical
facts.

The planning role maintains $G$ by first deriving an intermediate abstract
dependency graph from $\gamma$ and $\mathcal{F}$, then adding manipulator
assignment and execution dependencies. The final assignment-aware primitive
graph is
\begin{equation}
 \begin{aligned}
 G&=
 \rho_P\!\left(
 \gamma,\mathcal{F},\{\mathcal{C}_i\}_{i=1}^{N},\mathcal{R}\right)\\
 &\triangleq\bigl(\mathcal{N},\mathcal{L}\bigr),
 \end{aligned}
 \label{eq:planning_pipeline}
\end{equation}
where, for each $n\in\mathcal{N}$,
$n\triangleq\bigl(s_n,q_n,r_i,\kappa_n\bigr)$ denotes an assigned primitive
node with primitive skill type $s_n$, semantic target $q_n$, assigned
manipulator $r_i$, and execution skill $\kappa_n$. The relations in
$\mathcal{L}$ include the semantic task dependencies from the intermediate
dependency graph and manipulator-aware dependencies induced by $\mathcal{R}$.
Thus, $G$ records what to execute, who executes it, and which skill
is used. For $n\in\mathcal{N}$,
$\operatorname{pred}(n)$ denotes the predecessor nodes $j$ such that
$(j,n)\in\mathcal{L}$.

The verification role maintains $z$ after command execution. For an executable
command $c_n$, the command sent to the assigned manipulator interface specifies
the required execution skill and semantic target associated with primitive
subtask $n$. After $c_n$ is executed, the latest observation first refreshes
$\mathcal{F}$; the verification role then extracts transient execution
evidence
\begin{equation}
 \eta_n=\rho_V^{\texttt{e}}\bigl(c_n,\mathcal{F}\bigr),
 \label{eq:verification_evidence}
\end{equation}
where $\eta_n$ is the transient execution evidence for the executed command.
The verification role compares evidence with the current instruction to update
status and detect task-relevant events:
\begin{equation}
 \bigl(e_k,z(n)\bigr)=
 \rho_V^{\texttt{z}}\bigl(\eta_n,\gamma\bigr),
 \label{eq:verification_status}
\end{equation}
Here, $z(n)$ is the primitive-subtask status and $e_k$ records any detected
event in~\eqref{eq:semantic_events}; its affected scope is identified in
Section~\ref{sec:semantic_physical_coordination}. Without a task-relevant
inconsistency, $e_k$ is null and only $z(n)$ is updated.

Within $S_\nu^{\texttt{sem}}=(\gamma,G,\mathcal{F},z)$, the affordance,
planning, and verification roles own $\mathcal{F}$, $G$, and $z$, respectively.
An updated $\gamma$ enters as an external semantic constraint and is committed
before role-specific updates. The intermediate graph is used
only to construct $G$, and $\eta_n$ remains transient verification evidence.
